\documentclass[12pt]{article}
\usepackage[a4paper,margin=1in]{geometry}
\usepackage[utf8]{inputenc}
\usepackage[T1]{fontenc}
\usepackage{amsmath,amssymb}
\usepackage{textcomp}
\usepackage{parskip}
\begin{document}
\begin{center}
{\Large \textbf{2001 STEP I}}
\end{center}

\subsection*{Page 1}
\begin{flushleft}
\bigskip\textbf{Section A: Pure Mathematics}\par
\bigskip\textbf{Question 1.} The points A, B and C lie on the sides of a square of side 1 cm and no two points lie on the\par
same side. Show that the length of at least one side of the triangle ABC must be less than or\par
\[equal to (\sqrt\{6\} \sqrt\{2\}) cm.\]
-\par
\bigskip\textbf{Question 2.} Solve the inequalities\par
\textbf{(i)} 1 + 2x x2 > 2/x (x = 0) ,\par
- 6\par
\textbf{(ii)} \sqrt(3x + 10) > 2 + \sqrt(x + 4) (x > 10/3) .\par
-\par
3 Sketch, without calculating the stationary points, the graph of the function f (x) given by\par
\[f (x) = (x p)(x q)(x r) ,\]
- - -\par
\[where p < q < r. By considering the quadratic equation f (x) = 0, or otherwise, show that\]
\[(p + q + r) > 3(qr + rp + pq) .\]
By considering (x2 + gx + h)(x k), or otherwise, show that g2 > 4h is a sufficient condition\par
-\par
but not a necessary condition for the inequality\par
\[(g k) > 3(h gk)\]
- -\par
to hold.\par
\[3 tan θ tan3 θ\]
\bigskip\textbf{Question 4.} Show that tan 3θ = - .\par
\[1 3 tan2 θ\]
-\par
\[Given that θ = cos\]
\[-1(2/\sqrt\{5\})\]
\[and 0 < θ < \pi/2, show that tan 3θ = 11/2 .\]
Hence, or otherwise, find all solutions of the equations\par
-1\par
\textbf{(i)} tan(3 cos x) = 11/2 ,\par
1 -1\par
\textbf{(ii)} cos( tan y) = 2/\sqrt5 .\par
\end{flushleft}
\subsection*{Page 2}
\begin{flushleft}
\bigskip\textbf{Question 5.} Show that (for t > 0)\par
1 1\par
\textbf{(i)} dx = ,\par
\[Z^0 (1 + tx)2 (1 + t)\]
2x 1\par
\textbf{(ii)} - dx = .\par
\[Z^0 (1 + tx)3 - (1 + t)2\]
Noting that the right hand side of (ii) is the derivative of the right hand side of (i), conjecture\par
the value of\par
\[1 6x2\]
dx .\par
\[Z^0 (1 + x)4\]
(You need not verify your conjecture.)\par
\bigskip\textbf{Question 6.} A spherical loaf of bread is cut into parallel slices of equal thickness. Show that, after any\par
number of the slices have been eaten, the area of crust remaining is proportional to the number\par
of slices remaining.\par
A European ruling decrees that a parallel-sliced spherical loaf can only be referred to as 'crusty'\par
if the ratio of volume V (in cubic metres) of bread remaining to area A (in square metres) of\par
crust remaining after any number of slices have been eaten satisfies V /A < 1. Show that the\par
radius of a crusty parallel-sliced spherical loaf must be less than 2 metres.\par
[The area A and volume V formed by rotating a curve in the x-y plane round the x-axis from\par
\[x = a to x = a + t are given by\]
- -\par
-a+t\par
dy 2\par
-a+t\par
\[A = 2\pi y 1 + dx , V = \pi y dx . ]\]
Z-a \lambda \gamma dx \delta \mu Z-a\par
\end{flushleft}
\subsection*{Page 3}
\begin{flushleft}
\bigskip\textbf{Question 7.} In a cosmological model, the radius R of the universe is a function of the age t of the universe.\par
The function R satisfies the three conditions:\par
0 00\par
\[R(0) = 0, R (t) > 0 for t > 0, R (t) < 0 for t > 0, ( )\]
∗\par
where R denotes the second derivative of R. The function H is defined by\par
R (t)\par
\[H(t) = .\]
R(t)\par
\textbf{(i)} Sketch a graph of R(t). By considering a tangent to the graph, show that t < 1/H(t).\par
\textbf{(ii)} Observations reveal that H(t) = a/t, where a is constant. Derive an expression\par
for R(t). What range of values of a is consistent with the three conditions ( )?\par
∗\par
\textbf{(iii)} Suppose, instead, that observations reveal that H(t) = bt\par
-2,\par
where b is constant.\par
Show that this is not consistent with conditions ( ) for any value of b.\par
∗\par
\bigskip\textbf{Question 8.} Given that y = x and y = 1 x2 satisfy the differential equation\par
-\par
\[d2y dy\]
\[+ p (x) + q (x)y = 0 , ( )\]
\[dx2 dx ∗\]
\[show that p (x) = 2x(1 + x^2) -1 and q (x) = 2(1 + x^2) -1.\]
-\par
\[Show also that ax + b(1 x^2) satisfies the differential equation for any constants a and b.\]
-\par
Given instead that y = cos2(x2/2) and y = sin 2 (x2/2) satisfy the equation ( ), find p (x)\par
∗\par
and q (x).\par
\end{flushleft}
\subsection*{Page 4}
\begin{flushleft}
\bigskip\textbf{Section B: Mechanics}\par
\bigskip\textbf{Question 9.} A ship sails at 20 kilometres/hour in a straight line which is, at its closest, 1 kilometre from\par
a port. A tug-boat with maximum speed 12 kilometres/hour leaves the port and intercepts\par
the ship, leaving the port at the latest possible time for which the interception is still possible.\par
How far does the tug-boat travel?\par
\bigskip\textbf{Question 10.} A gun is sited on a horizontal plain and can fire shells in any direction and at any elevation\par
at speed v. The gun is a distance d from a straight railway line which crosses the plain,\par
where v2 > gd. The gunner aims to hit the line, choosing the direction and elevation so as to\par
maximize the time of flight of the shell. Show that the time of flight, T , of the shell satisfies\par
\[g 2 T 2 = 2v 2 + 2 v 4 g 2 d 2 2 .\]
-\par
\bigskip\textbf{Question 11.} A smooth cylinder with circular cross-section of radius a is held with its axis horizontal.\par
A light elastic band of unstretched length 2\pia and modulus of elasticity λ is wrapped round\par
the circumference of the cylinder, so that it forms a circle in a plane perpendicular to the axis\par
of the cylinder. A particle of mass m is then attached to the rubber band at its lowest point\par
and released from rest.\par
\textbf{(i)} Given that the particle falls to a distance 2a below the below the axis of the cylinder,\par
but no further, show that\par
\[9\pimg\]
\[λ = .\]
\[(3\sqrt\{3\} \pi)2\]
-\par
\textbf{(ii)} Given instead that the particle reaches its maximum speed at a distance 2a below\par
the axis of the cylinder, find a similar expression for λ .\par
\end{flushleft}
\subsection*{Page 5}
\begin{flushleft}
\bigskip\textbf{Section C: Probability and Statistics}\par
12 Four students, Arthur, Bertha, Chandra and Delilah, exchange gossip. When Arthur hears\par
a rumour, he tells it to one of the other three without saying who told it to him. He decides\par
whom to tell by choosing at random amongst the other three, omitting the ones that he knows\par
have already heard the rumour. When Bertha, Chandra or Delilah hear a rumour, they behave\par
in exactly the same way (even if they have already heard it themselves). The rumour stops\par
being passed round when it is heard by a student who knows that the other three already\par
have aready heard it.\par
Arthur starts a rumour and tells it to Chandra. By means of a tree diagram, or otherwise,\par
\[show that the probability that Arthur rehears it is 3/4.\]
Find also the probability that Bertha hears it twice and the probability that Chandra hears\par
it twice.\par
13 Four students, one of whom is a mathematician, take turns at washing up over a long period\par
of time. The number of plates broken by any student in this time obeys a Poisson distribution,\par
-λ n\par
the probability of any given student breaking n plates being e λ /n! for some fixed constant\par
λ, independent of the number of breakages by other students. Given that five plates are\par
broken, find the probability that three or more were broken by the mathematician.\par
\bigskip\textbf{Question 14.} On the basis of an interview, the N candidates for admission to a college are ranked in order\par
according to their mathematical potential. The candidates are interviewed in random order\par
(that is, each possible order is equally likely).\par
\textbf{(i)} Find the probability that the best amongst the first n candidates interviewed is the\par
best overall.\par
\textbf{(ii)} Find the probability that the best amongst the first n candidates interviewed is the\par
best or second best overall.\par
\[Verify your answers for the case N = 4, n = 2 by listing the possibilities.\]
\end{flushleft}
\end{document}
